\chapter{Conclusions and Future Work}
\label{chap:conclusion}

\section{Conclusion and Key Findings}
This thesis addressed a practical problem: clinical practice guidelines
from different organisations are distributed across heterogeneous documents
and are not directly available in a common representation for
cross-guideline querying and comparison. The approach taken was to construct
a knowledge graph that represents content from multiple guidelines in a
shared structure and then use that graph as a controlled source for
question generation.

The four-stage pipeline: ingestion, extraction, normalisation, and question
generation, was designed to handle the structural and terminological
variation that arises when documents from different publishing
organisations are processed together. The ingestion stage classified
document content into typed element records while preserving page-level
provenance. The extraction architecture separates entity types between a
rule-based layer and a contextual layer, with the full design supporting
LLM-based extraction for categories that require contextual interpretation.
The normalisation stage resolved duplicate entities, converted measurement
units to a common representation, and created equivalence edges between
matching targets from different organisations. The question-generation
stage selected graph edges as anchors, provided each anchor to the language
model as factual context, and checked the generated question for consistency
with that anchor before acceptance.

The pipeline was instantiated on eight documents from four organisations
covering Type~1 and Type~2 diabetes. The resulting knowledge graph contains
2,503 nodes across 13 node types and 3,485 edges across 18 relation types.
Normalisation created eight \textsc{equivalent\_to} edges linking matching
diagnostic thresholds across ADA, IDF, and Diabetes Canada after unit
conversion. These relations are not directly stated in any individual
source document; they arise from comparing independently stated targets
after normalisation.

The full GPT-4o-mini question-generation run produced 1,289 questions, of which 1,280 passed the rule-based validation checks and 9 were flagged by the rule-based validator because they contained numerical values absent from their anchors. A separate 491-question
catalogue was used for the detailed quality evaluation. From this catalogue, 30 questions were reviewed by two medical doctors:
26 were accepted without qualification, three contained clear
extraction-derived errors, and one was flagged for a clinical-scope
concern rather than a factual error.

In the reviewed sample, all three clear failures were traced to incorrect
graph relations rather than to the question-formulation step. This points to graph quality, particularly the correctness of extracted relations, as an important target for further improvement.

\section{Future Work}

Several directions follow from the work reported here.

The ingestion pipeline evaluated Docling as an alternative to PyMuPDF but did
not adopt it, primarily because of its runtime cost on the CPU-only development
hardware (Section~\ref{sec:exp_parsing}). A spot check on a small sample showed
Docling produced no ligature artefacts, unlike the PyMuPDF output, though the
comparison of word-spacing errors between the two was inconclusive. With GPU
acceleration available, re-evaluating Docling as the primary parser,
particularly for its ligature handling, would be a reasonable next step, and
could reduce or replace the family-specific text repairs currently required for
some document sources.

The LLM-based contextual extraction path was executed for the four entity categories of drugs, clinical conditions, populations, and interventions. However, extraction quality was not independently evaluated against a manually annotated gold-standard dataset, and the LLM-based extraction was not compared with the heuristic fallback mode of the same script.

Entity deduplication in this work is deterministic and relies on a fixed
synonym dictionary. Surface forms not represented in the dictionary remain
as separate nodes, which can increase the number of near-duplicate Drug and
ClinicalCondition nodes. Embedding-based or fuzzy matching could improve
recall in entity resolution, although safeguards would be needed to avoid
incorrectly merging distinct entities.

The current \textsc{equivalent\_to} criterion is implemented for normalised
plasma-glucose targets. Although alternate HbA1c representations are stored
during extraction, the equivalence procedure does not currently compare
HbA1c targets across organisations. Extending the criterion to HbA1c would
enable additional cross-guideline comparisons, including comparisons
involving NICE targets represented in the corpus.

The validator checks consistency between a generated question and its graph
anchor, but it does not determine whether the underlying graph relation is
clinically correct. A more complete evaluation would include a larger
reviewed sample, reviewers with a broader range of relevant clinical
expertise, and an analysis of whether the graph errors observed in the
current review are systematic or isolated.

Finally, the four-stage architecture is not tied to diabetes, but the
schema, extraction rules, canonical entities, repair procedures, and
normalisation configuration used here were developed for the diabetes
corpus. Applying the architecture to another multi-guideline domain would
test which components transfer with limited adaptation and which require
domain-specific redesign.